\section{D5 odd-$m$ final globalization theorem}

\subsection{Imported structural inputs}

We use the following imported propositions.

\begin{enumerate}
\item Explicit trigger family:
\[
H_{L1}=\{(q,w,u,\lambda)=(m-1,m-1,u,2):u\neq 2\}.
\]

\item Universal first-exit targets:
the regular families first exit at
\[
T_{\mathrm{reg}}=(m-1,m-2,1),
\]
while the exceptional family first exits at
\[
T_{\mathrm{exc}}=(m-2,m-1,1).
\]
\end{enumerate}

\subsection{Concrete bridge package}

On the $\Theta=2$ boundary section define
\[
\sigma \equiv w+u-q-1 \pmod m,
\qquad
\delta := q + m\sigma.
\]

On each splice-connected accessible component:
\[
\delta'=\delta+1 \pmod{m^2},
\]
and the current event class $\epsilon_4$ is a function of $(\beta,\delta)$.

\subsection{Globalization reductions}

\begin{enumerate}
\item Fixed realized $\delta$ can differ only by remaining tail length.
\item Every regular realization of every realized $\delta$ continues.
\item The exceptional continuation is pinned at chart/interface level to
\[
3m-3 \to 3m-2 \to 3m-1.
\]
\item Therefore the final globalization problem is the exceptional-row gluing
statement.
\end{enumerate}

\subsection{Final theorem}

\begin{theorem}
For odd $m$ in the accepted D5 regime, every actual lift of the exceptional
cutoff row $\delta=3m-3$ continues through
\[
3m-2 \to 3m-1
\]
and lies in the regular continuing endpoint class there.
Consequently, fixed realized $\delta$ determines endpoint class and future
word, hence
\[
\rho=\rho(\delta)
\]
globally, and raw global $(\beta,\delta)$ is exact.
\end{theorem}

\begin{proof}[Proof sketch]
The regular sector is already closed, so only the exceptional cutoff can be
terminal. The structural first-exit theorem gives a unique actual exceptional
exit target and branch direction. The interface theorem pins the corresponding
chain-label continuation to $3m-3 \to 3m-2 \to 3m-1$. Hence the exceptional
cutoff also has a forward successor. Therefore every accessible component is
total, every realization has infinite tail length, and the tail-length
reduction forces $\rho=\rho(\delta)$ globally.
\end{proof}

\subsection{Cleanup warning}

In a final manuscript, keep the following scope distinctions explicit:

\begin{itemize}
\item the concrete bridge package is componentwise before globalization;
\item the exceptional interface theorem is chart/interface level before it is
combined with the structural first-exit theorem;
\item the global theorem is obtained only after these inputs are combined.
\end{itemize}
